\section{System Architecture — Instantiation of the Formal Game}

The formal game defined in the previous section \ref{sec:formaldefinition} is realized in \package as a principled causal experimentation protocol — not as a loose software project.  
Its architecture is explicitly designed to preserve (i) the causal access restrictions, (ii) perfect reproducibility, and (iii) agent neutrality, while remaining fully modular and extensible.
These goals emerged from repeated attempts to deploy causal benchmarks whose behavior depended on hidden configuration or bespoke scripts \citep{RealCause2021,EvaluationFramework2022}. \package codifies the full stack—from SCM construction to metric evaluation—so that any research group can rerun the same experiment without oral tradition or manual tuning.

\subsection{Problem Instance as the Root Abstraction}
In implementation, a \emph{problem instance} is the root-level object that binds all components of a game together.  
It encapsulates exactly the tuple defined previously at the formal level:
\[
    \langle 
        \scm,\;\latentx \sqcup \measuredx \sqcup \controllablex,\;
        \initstate,\; \phi,\;
        \goal,\; \resultmetric,\; \behaviormetric
    \rangle.
\]

However — crucially — this is constructed \emph{not in code directly}, but via a fully serializable \textbf{Spec configuration}.  
This ensures:

\begin{itemize}
    \item \textbf{Perfect reproducibility}: any experiment can be saved and replayed bit-accurately.
    \item \textbf{Pluggability}: replacing the SCM, mission, or agent requires no change to the runtime engine.
    \item \textbf{Parallel evaluation}: multiple instances can be run concurrently, with isolated seeds and transcripts.
\end{itemize}

\noindent
Every execution of \game begins by \emph{loading a problem instance from Spec}, then passing it to the \textbf{Runtime}, which produces a \textbf{Transcript} — the immutable record of the agent’s interaction with the SCM.
Because Specs are plain data, they can be version-controlled, shared, and peer-reviewed just like benchmark definitions for treatment-effect suites or graph challenges \citep{CauseBox2021,CausalBench2024}.

\medskip
What follows is a deep dive through each architectural layer — beginning with the role of Contracts, the foundation that makes \game not only a tool, but a \emph{scientific protocol}.

\subsection{Contracts — Enforcing Causal Boundaries}

Every core component in \game (agent, SCM, mission, metric, hook, plot, etc.) is defined first by an abstract \emph{contract}.  
A contract specifies \emph{what a component is allowed to do} — never \emph{how it must be implemented}.  
This mirrors scientific protocol design: an experimental instrument is defined by its \emph{accessible interface}, not by its internal mechanics.

The agent contract, for example, can act and be informed, but cannot inspect the SCM.  
The SCM contract can generate interventional data but knows nothing about the agent’s beliefs.  
Result metrics see the ground truth; agents never do.  
This disciplines the entire system to respect causal epistemic validity — knowledge must be earned, not granted.
By separating capability from implementation, \game enables community contributions without compromising scientific integrity: two teams can implement radically different agents while remaining confident that neither receives unintended privileges \citep{BayesianCausalDiscovery2023}.

\subsection{Specs — Declarative Instantiation of a Game Instance}
\label{sec:specs}

A game instance is not created manually in Python code.  
Instead, every component is declared via a \emph{Spec}: a fully serializable description of
which class to use (via import path) and with which parameter configuration.  
A centralized registry instantiates these Specs at runtime — deterministically and seed-bound.

This makes \game \emph{perfectly reproducible} and experiment definitions \emph{immutably declarative}:  
a problem instance is not code, but data.  
Replacing the agent or SCM requires no change to the runtime engine — only to the Spec.
This declarative approach extends ideas from causal modeling libraries such as DoWhy \citep{DoWhy2021} but pushes them to benchmark scale: every executable element of the experiment can be reconstructed from a single YAML or JSON payload.

\subsection{DTOs — Controlled Information Flow}

During execution, every piece of information passed between components flows solely through \emph{Data Transfer Objects (DTOs)} — strictly typed, immutable containers.  
They act as \emph{information firewalls}, ensuring that no component receives more data than its contract permits.

For example, an agent receives only observable samples and round metadata — never latent variables or SCM internals.  
Metrics and plots may later access the full transcript and ground truth — but only after the game has ended.

DTOs thus preserve causal fairness even when evaluating extremely powerful or self-reflective agents.
Their schemas, documented in Appendix~\ref{appendix:dtos}, have been engineered so that even large language model-based agents—prone to exploiting incidental metadata—cannot infer privileged information.

\subsection{Managers — Instrumentation and Enforcement}
Managers regulate the execution of \game without influencing its causal logic. They do not assist the agent or modify the SCM; instead, they enforce constraints and expose observability.

\paragraph{BudgetManager} charges resource costs (e.g.\ samples, rounds, time, memory) and triggers termination via~$\phi$ when limits are exceeded — making experimentation economically meaningful.

\paragraph{HookManager} listens to lifecycle events (e.g.\ \texttt{GAME\_START}, \texttt{ROUND\_END}, \texttt{GAME\_END}) and triggers user-defined side-channel observers such as loggers or monitors. Hooks are strictly passive: they can record or visualize, but never intervene.

\paragraph{PlotManager} generates visual artifacts post-hoc (e.g.\ at game or benchmark end), never visible to the agent during play, ensuring causal purity while allowing researchers to analyse transcripts visually without contaminating the interactive loop.

Together, these managers provide accountability, observability, and control — never assistance or information leakage. They make the latent constraints that govern real-world experimentation (budgets, safety audits, post-hoc reporting) explicit in software \citep{Prosperi2020CausalIAA,ACCBench2024}.

In short, managers provide accountability, observability, and control — never assistance or information leakage.

\subsection{SCM Construction Process}
\game allows for the construction of an explicit, fully specified causal model comprising a directed acyclic graph (DAG), semantically typed variables, and structural equations with exogenous noise. This process ensures that every resulting SCM is identifiable, controllable, and suitable for scientifically valid causal experimentation.

\subsubsection{DAG Topology and Structural Constraints}
The causal graph is defined as a directed acyclic graph
\[
    \DAG = (\varset, \edgesset), \quad \edgesset \subseteq \varset \times \varset,
\]
where an edge $(\varname_j \to \varname_i)$ implies direct causal influence via the structural equation of $\varname_i$.  
DAG generation is never arbitrary — it is constructed to enforce causal identifiability, appropriate difficulty, and benchmark realism.

A \emph{strict topological ordering} is sampled such that all edges respect $j < i$, guaranteeing acyclicity and enabling forward-sampling during data generation.  
The DAG is aligned with the access partition
\[
    \latentx \sqcup \measuredx \sqcup \controllablex = \varset,
\]
where latent variables may act as confounders, only measured variables are observable to the agent, and only controllable variables may be intervened on during gameplay. This partition is fixed once the SCM is instantiated and cannot change mid-run.

Causal complexity is explicitly controlled by bounding the in-degree of each node (e.g.\ $\deg^-(\varname_i) \le \dmax$) and the maximum causal depth. This prevents trivial graphs and also avoids pathological, overly deep causal chains.  
Unlike most existing benchmarks, \game natively supports \emph{hybrid} causal graphs containing both continuous and categorical nodes. Variable type proportions are user-controlled, enabling diverse families of causal structures.

Finally, while the DAG may encode arbitrarily rich parent sets, its structure is never leaked — the agent cannot inspect or query $\edgesset$ directly. All knowledge arises solely through interventions and observed consequences.

\subsubsection{Variable Definitions: Type, Domain, and Access}
Once the graph structure is fixed, each variable $\varname_i \in \varset$ is assigned a semantic specification determining how it participates in the SCM.

Each variable has a declared type
\[
    \text{type}(\varname_i) \in \{\text{categorical},\ \text{continuous}\},
\]
with proportions governed by a user-defined parameter $p_C$. Continuous variables operate on bounded real intervals, while categorical variables draw from a finite, user-configurable domain $\{c_1, \dots, c_K\}$. Domains are immutable — interventions may only assign values already in the domain, never extend it.

In parallel, each variable is assigned an access role:
\[
    \text{access}(\varname_i) \in \{\text{latent},\ \text{observable},\ \text{controllable}\}.
\]
Latent variables never appear in any dataset, observable variables appear in all samples returned to the agent, and controllable variables may be directly set through interventions. This access role is contractual and strictly enforced throughout the game.

At this point, each variable has a clearly defined type, legal value range, position in the DAG, and epistemic status relative to the agent — forming a complete “vocabulary” for the SCM.

\subsubsection{Structural Equations and Noise Generation}
The final layer of construction defines how each variable’s value is \emph{generated} during sampling. For each variable $\varname_i$ with parent set $\mathrm{Pa}(\varname_i)$, the SCM defines a structural assignment
\[
    \varname_i := f_i\big(\mathrm{Pa}(\varname_i)\big) + \eta_i,
\]
where $f_i$ is a deterministic function and $\eta_i$ is an exogenous noise term, independently drawn per variable.

In SCMs generated from scratch, $f_i$ is sampled from a symbolic function grammar with configurable expressiveness — supporting operators such as addition, multiplication, exponentiation, logarithm, trigonometric functions, and bounded polynomial degrees. The function is generated exactly once at SCM creation time and remains immutable thereafter.

For categorical targets, class-specific scores are generated and normalized into a valid discrete distribution.  
Noise variables $\eta_i$ are drawn from family-specific distributions (Gaussian, Laplace, uniform, Dirichlet, or degenerate Dirac for fixed treatment variables), with output values always clipped or transformed to remain in the declared domain of $\varname_i$.

This mechanism cleanly separates deterministic functional influence from stochastic exogenous uncertainty — allowing the SCM to express richly structured causal phenomena while remaining perfectly simulatable under interventions.

---

This completes the SCM construction process: a topologically valid, access-partitioned DAG equipped with semantically defined variables and precise structural equations — ready to serve as the causal ground truth for interactive experimentation in \game.

\subsection{Missions — Declarative Definition of the Causal Task}

A \textbf{Mission} is a declarative, contract-based specification that defines \emph{what causal problem the agent is expected to solve}. It encodes the semantic nature of the task — for example, learning the full DAG, estimating the ATE or CATE, or reconstructing structural mechanisms — but it does \textbf{not} evaluate correctness itself. Instead, a mission orchestrates \emph{which output format} the agent must produce ($\outputspace$) and \emph{which evaluators apply} (i.e., which \textbf{result} and \textbf{behavior} metrics will later be invoked).

In other words, missions do not compute scores — they define the rules of the scientific game. They specify the expected answer representation and the evaluative lenses through which agent performance will be interpreted. This makes them a purely semantic layer — strictly framing \emph{what “solving the task” means} without influencing the environment or agent logic.

Because missions are declared via Specs, switching from “full DAG discovery” to “ATE estimation” requires no modification to the agent or SCM — only a change in mission configuration. Missions therefore serve as the \emph{semantic lens} through which causal inference is framed and evaluated in \game, keeping the system modular and scientifically expressive.
Concrete mission archetypes and their evaluation hooks are detailed in Section~\ref{sec:missions_examples}.

\subsection{The Agent Perspective — Action, Observation, and Learning}
From the environment’s point of view, the SCM is the causal world. From the agent’s point of view, the world is experienced only as a sequence of DTO-restricted observations and permissible action opportunities. The agent operates entirely under partial observability, without privileged access to the true SCM, while navigating unknown ground truth and being subject to explicit resource constraints.

Every agent conforming to the \texttt{Agent} contract implements two methods: \texttt{act(...)} and \texttt{inform(...)}. At the beginning of each round, the agent is invoked with \texttt{act(RoundInformation, AvailableActions)}, where it receives only legal meta-information and the currently admissible action space. It must then propose either a valid experiment or a terminal \emph{stop with answer} action. After the environment executes that request, the agent is called with \texttt{inform(Samples, Feedback)}, receiving only the observable new samples and optional mission-defined feedback — never any internal representation of the SCM. Learning agents may use this information to update latent internal beliefs, but have no direct access to the transcript, budget counters, or ground truth.

This observational restriction is absolute. The agent never sees SCM equations $f_i$, noise variables $\eta_i$, latent variables $\latentx$, future reward signals, or the true causal structure. It operates solely on previously observed samples (restricted to $\measuredx$), optional round-level metadata provided through DTOs, and the explicitly declared permissible action space.

At any time, the agent may choose to stop voluntarily by submitting a final answer $y \in \outputspace$. This ends the game immediately. If it fails to do so before the stopping function $\phi$ triggers, termination occurs without an answer, and the result metric assigns a failure score.

Crucially, \game is architecturally neutral: it does not assume, privilege, or restrict agents to any particular paradigm. Agents may be symbolic search procedures, Bayesian experimental designers, deep learning models, reinforcement learners, or any hybrid thereof. All are subjected to the exact same causal interaction protocol — no agent is granted information outside legal DTOs, and no agent may bypass the constraints of the environment.

With the agent interface now fully stated, we proceed to define how \game evaluates performance — not only in terms of correctness of the final answer, but also in terms of the efficiency and prudence of the experimental process.

\subsection{Result Validators — Ensuring Output-Format Legitimacy}
A \textbf{Result Validator} is the gatekeeper that checks whether the agent’s submitted answer $y \in \outputspace$ is \emph{well-formed and admissible} — before any scoring is attempted. Its role is strictly \emph{pre-evaluative}: it ensures that the output matches the mission-defined schema, datatype, and structural constraints.

This validation step prevents malformed or logically invalid outputs (e.g.\ ill-shaped tensors, non-DAG matrices, out-of-domain predictions) from being passed to the \resultmetric. If the output is invalid, the validator can either reject it outright or return a canonical “invalid answer” token for scoring. Crucially — \textbf{the validator does not judge correctness} — it only ensures that the answer is structurally legitimate input for the result metric.

Once validation succeeds, the \textbf{result metric} takes over, comparing the validated output to the mission’s ground truth and assigning the actual correctness score.

\subsection{Evaluation Objectives — Behavioral vs.\ Result Metrics}
\label{sec:eval:metrics}
A defining principle of \game is that agents are \emph{not judged solely on correctness}. Their performance is evaluated along two orthogonal dimensions: the quality of the causal answer they ultimately submit, and the efficiency with which they conducted the experimental process required to obtain it.

The first dimension is the \textbf{result metric} $\resultmetric : \outputspace \rightarrow \mathbb{R}$, which evaluates the agent’s final answer $y$ after it voluntarily issues a stop-with-answer action. The score reflects how well $y$ matches the ground truth, and may take the form of Structural Hamming Distance or F1 for DAG recovery, PEHE or mean-squared error for treatment effect estimation, or functional similarity or KL divergence for full SCM reconstruction. If the agent fails to return an answer before the stopping function $\phi$ is triggered, it receives a mission-dependent minimal score (e.g.\ $0$ or $-\infty$).

The second dimension is the \textbf{behavior metric} $\behaviormetric : \runs \rightarrow \mathbb{R}$, which evaluates the entire trajectory $r$ rather than its final outcome. This reflects the agent’s experimental prudence — for example, by penalizing excessive sample usage, intervention count, or round duration via a task-specific cost model.

Crucially, \game does not combine these scores into a single scalar value. Instead, it returns the two-objective outcome
\[
    \big( \resultmetric(y),\ \behaviormetric(r) \big),
\]
allowing users to apply their preferred aggregation — such as convex combination, lexicographic ordering, or explicit Pareto frontier analysis.

By separating correctness from experimentation strategy, \game makes efficiency–accuracy trade-offs \textbf{formally measurable}. It distinguishes, for the first time at benchmark scale, between an agent that is perfectly accurate but recklessly expensive, and an agent that is slightly less accurate but vastly more efficient.

\subsection{Serialization and Artifacts}
\label{sec:serialization}
A defining property of \game is that every execution is \textbf{perfectly reproducible}, \textbf{auditable}, and \textbf{replayable} — down to the exact sequence of interventions chosen, samples observed, and decisions made. This is made possible through strict serialization at the level of experimental definition and explicit artifact generation at runtime.

Because every game instance is declared entirely via Specs (cf.\ Sec.~\ref{sec:specs}) and every SCM is executed under a fixed global seed, an experiment is serializable \emph{before it even begins}. The full configuration — including SCM provenance, agent parameters and seeds, mission definitions, metric specifications, and budget or hook rules — is recorded deterministically.

At runtime, the system produces a \textbf{Transcript artifact}: a structured, DTO-encoded log that captures every intervention request, every sample revealed to the agent, every decision made (including the final answer, if any), and every applied cost or forced stopping event. This transcript is immutable, machine-parseable, human-inspectable, and replayable — it is the single source of truth for evaluation, visualization, and scientific auditing.

This architecture ensures that any experiment conducted in \game can be reproduced exactly across machines, audited rigorously for fairness or correctness, revisited with new evaluation tools, and safely shared as a benchmark asset — all without leaking ground-truth causal information prematurely.

With serialization guaranteed, we now turn to how \game scales this rigor to large-scale empirical evaluation — enabling the concurrent benchmarking of many agents, SCMs, or missions in parallel.

\subsection{Runner and Parallel Execution}
\label{sec:runner}
While a single \textbf{Game} represents one agent–SCM interaction, most real experimental scenarios require running many such games — for example, comparing agents across different SCMs, stress-testing a single agent over many seeds, or sweeping entire agent × SCM × mission grids. To support this, \game provides a dedicated orchestration layer called the \textbf{Runner}.

The Runner automates the execution of multiple game instances, either sequentially or in parallel, while guaranteeing full isolation: each run receives its own SCM instance (even if derived from the same Spec), its own agent and random seed, and its own transcript. No run may observe or influence another — preserving experimental purity even in large-scale or distributed evaluation settings.

Once execution finishes, the Runner performs systematic post-processing. It surfaces a unified collection of transcripts, aggregated metric evaluations, and — where specified — benchmark-level visualizations or Pareto frontier analyses triggered by the PlotManager. In doing so, it does not modify or reinterpret the underlying causal game; it merely manages repetition, isolation, and aggregation.

\subsection{A Scientific Protocol, Not Just Software}

The architecture above is deliberately minimal, but precisely controlled.  
Every component is replaceable, but only through declared Specs.  
Every piece of information is visible only where epistemically valid.  
Every run is fully replayable.  
\game is therefore not a heuristic system, but a \textbf{formal causal experimentation protocol instantiated in software}.

\medskip
With the architecture now formalized as an \emph{executable causal protocol}, we are ready to define how it is used — the \emph{experimental protocol} that governs evaluation, comparability, and scientific validity.
